\documentclass[11pt]{article}

\usepackage[margin=1in]{geometry}
\usepackage{booktabs}
\usepackage{hyperref}
\usepackage{amsmath}
\usepackage{amssymb}

\title{Conservative Generalization over Assumption Graphs:\\
A Bounded Framework-Evolution Agent for Recursive Self-Evolution}
\author{Anonymous}
\date{}

\begin{document}
\maketitle

\begin{abstract}
We study whether an agent can improve by proposing new assumptions, validating
them against old successes and residual failures, and selectively retaining only
bounded generalizations.  The system represents assumptions as an Assumption
Graph and promotes local repairs into candidate frameworks only when they pass
Conservative Generalization obligations: old-success preservation, residual
explanation, limiting-case reduction, new prediction, and regression control.
The result is a bounded L3.5 recursive self-evolution prototype, not an
unbounded autonomous OS.
\end{abstract}

\section{Introduction}

The central question is not whether a model can produce another prompt.  The
question is whether a system can turn failure into a structured hypothesis,
validate whether the hypothesis preserves prior competence, and retain it as a
new framework only when it explains more than a local patch.  We call this
process bounded dialectical framework growth.

\section{Related Work}

The system is positioned relative to RAG and graph memory, Reflexion-style
self-improvement, evolutionary agents, world-model-guided planning, and formal
analogy systems.  The core distinction is that this work evaluates the
framework-growth step itself: candidate frameworks must reduce to parent
frameworks under limiting cases and must predict new validation rows.

\section{Assumption Graph}

Each hypothesis is represented as a typed node with evidence, trials, residual
labels, verifier obligations, and lifecycle status.  Edges encode relations
such as generalizes, reduces-to-under-scope, explains-residual,
preserves-success-cases, modifies-boundary-of, predicts-new-case, conflicts-with,
and orthogonal-to.  This graph is the memory substrate for recursive
self-evolution.

\section{Dialectical Framework Growth}

The agent separates three levels of new assumptions:

\begin{itemize}
  \item branch: a scoped local repair or negative-evidence branch;
  \item candidate framework: a structured generalization with full validation obligations;
  \item active scoped framework: a retained framework with bounded scope, certificates, and recheck requirements.
\end{itemize}

The Hegelian structure is operational rather than rhetorical.  A parent
framework is the thesis, a residual cluster is the antithesis, and the candidate
framework is the synthesis.  Selective retention requires that the synthesis
preserve old successes while explaining residuals and producing new predictions.

\section{Conservative Generalization}

For a candidate framework \(F_{\mathrm{new}}\) and parent frameworks
\(F_{\mathrm{old}}\), the framework growth score is computed from:

\[
\mathrm{Growth}(F_{\mathrm{new}} | F_{\mathrm{old}}) =
P_{\mathrm{old}} + E_{\mathrm{residual}} + R_{\mathrm{limit}} +
G_{\mathrm{scope}} + N_{\mathrm{prediction}} - C_{\mathrm{regression}}.
\]

Promotion is blocked when old-success preservation fails, controls are harmed,
or the candidate cannot state limiting-case reductions.  This prevents raw
wisdom strings and longer prompts from being mistaken for framework growth.

\section{Simulator-Guided Verification}

The world model is used as a cheap router and budget gate.  It selects verifier
tiers and fresh-validation candidates, but it is not allowed to replace live
ablation or judge calls.  Simulator defects are written back as residuals for the
next generation.

\section{Finite Formal Certificates}

The formal layer is category-inspired and bounded.  It checks finite structural
morphism certificates, invariant preservation, negative controls, and Lean-
checked theorem fragments where applicable.  It does not claim to be a full
category-theory theorem prover.

\section{Experiments}

\subsection{Frozen Main Line}

Report the frozen benchmark line: tasks, baselines, problem-level bootstrap
confidence intervals, and domain breakdown.

\subsection{Fresh Repaired Generator Rerun}

Report the 720-call repaired broad-generator rerun and the distinction between
gated selective retention and the failed unfiltered generator claim.

\subsection{Framework Growth Benchmark}

Compare the full framework-evolution agent against local patching, raw wisdom,
no morphism, no world model, no lifecycle ledger, and no conservative gate.

\subsection{LLM Candidate Framework Experiment}

Report the residual-driven LLM-framework-candidate experiment: ten candidate
frameworks, top-two conservative validation, negative control retention, and the
claim boundary for live LLM API generation.

\subsection{External Evaluation Pack}

Report the expert annotation packet, proxy preflight, artifact hash manifest,
redaction policy, and exact reproducibility commands.  Do not claim a completed
human panel until it has been run.

\section{Negative Results and Claim Boundaries}

The system cannot claim unbounded autonomous self-evolution, a production world
simulator replacing live validation, or a full category-theory theorem prover.
Broad unfiltered generation failed before repaired selection.  Some gains depend
on gating and abstention rather than universal activation.

\section{Reproducibility}

The release includes machine-readable artifacts, exact commands, redacted
payloads, claim ledgers, and coverage audits.  API credentials must be supplied
through environment variables and must not be committed to code or artifacts.

\section{Conclusion}

The contribution is a bounded framework-evolution agent: it can propose,
validate, retain, demote, and audit new assumption frameworks over an Assumption
Graph.  The result supports a paper claim of conservative generalization and
recursive self-evolution under explicit gates, not a claim of unbounded autonomy.

\end{document}
